\scanpage{52}
\mspage{43}
% archive sheet number: 52

Fixons un $i\in I$. Pour tout $j\leq i$, posons
\[
\begin{aligned}
X_j^{\circ\prime}&=\text{cellule \unsure{ouverte} de }X_j\text{ dans }\unsure{\Phi|P}_i\\
&=X_j\setminus\bigcup_{\substack{j'\leq i\\\text{t.q. }X_{j'}\subsetneq X_j}}X_{j'}
\end{aligned}
\]
\margin{\uline{NB} dans le lemme, on ne suppose \uline{pas} $j'\leq j$}

\inserted{Je dis que} cette somme comporte plus de termes a priori que la somme analogue qui figure dans la déf. de
\[
X_j^\circ=X_j\setminus\bigcup_{j'<j}X_{j'},
\]
i.e. que
\[
j'<j\Longrightarrow X_{j'}\subsetneq X_j.
\]
En effet, on \deleted{\unsure{sait}} \inserted{a} \deleted{que $X_j$} $X_{j'}\subset\partial X_j$, donc \deleted{\unsure{on}} donc $X_j\setminus X_{j'}\supset X_j\setminus\partial X_j=X_j^\circ$, comme $X_j^\circ\ne\varnothing$ par hyp. on conclut \unsure{que} bien $X_j\setminus X_{j'}\ne\varnothing$ i.e. $X_{j'}\ne X_j$.

On a donc prouvé par là-même
\[
X_j^{\circ\prime}\subset X_j^\circ\qquad\text{``\unsure{bigre|bigger}''}
\]

\deleted{d'autre part par hyp. Mais l'ens des $X_j^{\circ\prime}$ \unsure{recouvre} $X_i$} \inserted{(\unsure{par la publ.} :)} \deleted{Or \unsure{a}} \deleted{Donc}
\[
\deleted{\bigcup_{j\leq i}X_j^{\circ\prime}\subset\bigcup_{j\leq i}X_j^\circ\subset X_i}
\]
% slanted
\margin{donc $\bigcup X_j^{\circ\prime}=X_i$ implique $\bigcup X_j^\circ=X_i$, qfd. \unsure{$\dagger$}}

\deleted{avec \uline{égalité} des deux extrêmes. Comme les $X_j^\circ$ sont $\ne\varnothing$ et deux à deux disjoints, cela implique}
\[
\deleted{X_j^{\circ\prime}=X_j^\circ:}
\]
\deleted{Les cellules ouvertes sont les $X_j^\circ$. En même} \deleted{En même temps,} \deleted{cela entraîne que} \deleted{temps, si $j,k\leq i$, je dis que}
\[
\deleted{\unsure{i}}\qquad\deleted{j\unsure{\nleq|\ne}k\Longrightarrow}
\]

\deleted{Soient $i,j\in I$ tels que $X_i\subset X_j$, prouvons $i\leq j$. Soit $x\in X_i^\circ$ ($X_i^\circ\ne\varnothing$ par $b')$}
